% =========================================================
% 6.6 DISEÑO DE INTERFACES DE USUARIO (UI/UX)
% El \section{} de esta sección está en el chapter.tex del capítulo;
% sus referencias van en seccion6_6.bib (misma carpeta).
% Las pantallas se exportan de Figma al doble de resolución (720 x 1600) y
% se convierten a JPEG de 540 x 1200 en
% figures/capitulo6_diseno_sistema/6.6_Diseno_de_interfaces_de_usuario/,
% subcarpetas modo_claro/ y modo_oscuro/, con el mismo nombre en ambas.
% El prefijo indica el flujo: sesion_, mapas_, rutas_, navegacion_,
% reportes_, favoritos_, monitoreo_ y perfil_. Cada figura agrupa hasta seis
% pantallas como subfiguras (macro \pantallaUI, setup/settings.tex, sección E).
% =========================================================
Las interfaces de usuario del prototipo se diseñaron en Figma como maquetas de alta fidelidad para un dispositivo móvil con una resolución de $360 \times 800$ píxeles. En esta sección se presentan los elementos de identidad visual de la aplicación, es decir, el logo y la paleta de colores, así como las pantallas que la conforman en sus modos claro y oscuro.

\subsection{Logo}
El logo de la aplicación es un perro (figura \ref{fig:logo}). Se eligió esta figura porque el perro es un símbolo ampliamente asociado con la lealtad, la compañía y la protección, cualidades que la aplicación busca transmitir al acompañar al peatón en sus trayectos. Con ello se pretende generar confianza en los usuarios, aspecto fundamental en una herramienta cuyo propósito es sugerir rutas seguras, así como lograr una imagen de impacto que sea fácil de reconocer y recordar.

\begin{figure}[htbp]
    \centering
    \includegraphics[width=0.5\linewidth]{figures/capitulo6_diseno_sistema/6.6_Diseno_de_interfaces_de_usuario/logo.png}
    \caption{Logo de la aplicación}
    \label{fig:logo}
\end{figure}

\subsection{Paleta de colores}
La paleta de colores de la aplicación parte de tres tonos principales: naranja terracota, azul y verde. En conjunto, estos tonos hacen alusión a un atardecer: el azul remite al cielo, el naranja terracota a la luz del sol al ocultarse y el verde a la vegetación del paisaje, lo que da a la interfaz una apariencia cálida y armónica.

Los colores se definieron en Figma como variables agrupadas en colecciones según su función. Para el modo oscuro se definió una colección adicional con las mismas variables y valores adaptados a fondos oscuros, de modo que ambos modos comparten la misma estructura y solo cambian los valores. El cuadro \ref{tab:paleta_colores} muestra los valores de cada variable en ambos modos, y su uso dentro de la interfaz es el siguiente:

\begin{itemize}
    \item \textbf{Marca.} El color primario (\texttt{primary}), naranja terracota, es el color principal de la aplicación: se emplea en los botones de acción, en el elemento activo de la barra de navegación, en el botón para añadir lugares y en el marcador de ubicación sobre el mapa. El color secundario (\texttt{secondary}), azul, funciona como color de apoyo, por ejemplo en los campos del formulario de edición de perfil.
    \item \textbf{Semáforo.} El verde, el amarillo y el rojo comunican el estado de una acción: el verde se usa en los mensajes de confirmación y en el trazo de la ruta; el amarillo, en advertencias, como el aviso de límite de reportes, y el rojo, en acciones destructivas, como la eliminación de la cuenta.
    \item \textbf{Acento.} Son versiones suaves de los colores de marca y de semáforo que se usan como fondo de componentes que deben destacar sin saturar la pantalla, como el encabezado de las pantallas de inicio de sesión y perfil, la categoría seleccionada al reportar un incidente y los mensajes de confirmación, advertencia y eliminación.
    \item \textbf{Fondos.} Definen el fondo principal de las pantallas y un fondo secundario para elementos como los campos de texto. En el modo oscuro se añade una superficie intermedia (\texttt{surface}) para distinguir los contenedores del fondo.
    \item \textbf{Texto.} Distinguen el texto principal del texto que se coloca sobre fondos de color, como el de los botones de acción.
\end{itemize}

\begin{table}[htbp]
    \centering
    \caption{Paleta de colores de la aplicación en los modos claro y oscuro}
    \label{tab:paleta_colores}
    \begin{tabular}{|l|l|l|}
        \hline
        \textbf{Variable} & \textbf{Modo claro} & \textbf{Modo oscuro} \\ \hline
        \multicolumn{3}{|l|}{\textit{Marca}} \\ \hline
        \texttt{primary}                    & \muestraColor{C85A17} & \muestraColor{D9661F} \\ \hline
        \texttt{secondary}                  & \muestraColor{BFE3FA} & \muestraColor{3B6A8A} \\ \hline
        \multicolumn{3}{|l|}{\textit{Semáforo}} \\ \hline
        \texttt{green}                      & \muestraColor{2E7D4F} & \muestraColor{3DA66A} \\ \hline
        \texttt{yellow}                     & \muestraColor{E0A800} & \muestraColor{F0B90B} \\ \hline
        \texttt{red}                        & \muestraColor{D0212D} & \muestraColor{EF4B55} \\ \hline
        \multicolumn{3}{|l|}{\textit{Acento}} \\ \hline
        \texttt{primary\_accent}            & \muestraColor{FBE4D6} & \muestraColor{3A2418} \\ \hline
        \texttt{secondary\_accent}          & \muestraColor{E9F7FF} & \muestraColor{16252E} \\ \hline
        \texttt{green\_accent}              & \muestraColor{E1F9EB} & \muestraColor{143323} \\ \hline
        \texttt{yellow\_accent}             & \muestraColor{FEF2CD} & \muestraColor{3A3014} \\ \hline
        \texttt{red\_accent}                & \muestraColor{F5D8DA} & \muestraColor{3F1D21} \\ \hline
        \multicolumn{3}{|l|}{\textit{Fondos}} \\ \hline
        \texttt{background\_white}           & \muestraColor{FFFFFF} & \muestraColor{121212} \\ \hline
        \texttt{background\_black}           & \muestraColor{000000} & \muestraColor{F2F2F2} \\ \hline
        \texttt{background\_white\_secondary} & \muestraColor{F4F4F4} & \muestraColor{2A2A2A} \\ \hline
        \texttt{background\_black\_secondary} & \muestraColor{9D9D9D} & \muestraColor{A3A3A3} \\ \hline
        \texttt{surface}                    & ---                   & \muestraColor{1E1E1E} \\ \hline
        \multicolumn{3}{|l|}{\textit{Texto}} \\ \hline
        \texttt{Black}                      & \muestraColor{000000} & \muestraColor{F2F2F2} \\ \hline
        \texttt{White}                      & \muestraColor{FFFFFF} & \muestraColor{FFFFFF} \\ \hline
    \end{tabular}
\end{table}

En el modo oscuro, el color primario y los colores de semáforo se aclaran ligeramente para conservar su contraste sobre fondos oscuros; por ejemplo, el primario pasa de \texttt{\#C85A17} a \texttt{\#D9661F}. El secundario, en cambio, pasa de un azul claro a uno más profundo (\texttt{\#3B6A8A}), y los acentos se convierten en tonos oscuros, como el café del encabezado (\texttt{\#3A2418}). Además, el fondo principal no es negro puro sino un gris muy oscuro (\texttt{\#121212}) y el texto principal se invierte a un blanco grisáceo (\texttt{\#F2F2F2}), lo que evita el contraste extremo entre el texto y el fondo.

\subsection{Modo claro}
La aplicación se organiza en cuatro secciones, accesibles desde la barra de navegación inferior: Perfil, Explora (el mapa), Favoritos y Monitoreo. Las figuras \ref{fig:ui_sesion_registro} a \ref{fig:ui_perfil_cuenta} muestran las 94 pantallas del modo claro agrupadas por flujo, y la tabla \ref{tab:ui_flujos} relaciona cada flujo con sus figuras y con los requerimientos funcionales que atiende.

\begin{table}[htbp]
    \centering
    \caption{Flujos de pantallas y requerimientos que atienden}
    \label{tab:ui_flujos}
    \begin{tabular}{|>{\raggedright\arraybackslash}p{0.38\textwidth}|>{\raggedright\arraybackslash}p{0.22\textwidth}|>{\raggedright\arraybackslash}p{0.27\textwidth}|}
        \hline
        \textbf{Flujo} & \textbf{Figuras} & \textbf{Requerimientos} \\ \hline
        Sesión, registro y recuperación de contraseña & \ref{fig:ui_sesion_registro} a \ref{fig:ui_sesion_recuperacion} & RF-01 a RF-07, RF-10, RF-11, RF-13 \\ \hline
        Mapa y búsqueda de lugares & \ref{fig:ui_mapa_busqueda} y \ref{fig:ui_mapa_lugar} & RF-15 a RF-17, RF-30 \\ \hline
        Cálculo y selección de rutas & \ref{fig:ui_rutas} y \ref{fig:ui_rutas_franja} & RF-22 a RF-26, RF-28 \\ \hline
        Navegación & \ref{fig:ui_navegacion} y \ref{fig:ui_navegacion_fin} & RF-27, RF-29 a RF-35, RF-47 \\ \hline
        Reportes & \ref{fig:ui_reportes_registro} a \ref{fig:ui_reportes_resultado} & RF-20, RF-37 a RF-42 \\ \hline
        Lugares favoritos & \ref{fig:ui_favoritos} y \ref{fig:ui_favoritos_edicion} & RF-18 \\ \hline
        Monitoreo & \ref{fig:ui_monitoreo_enlaces} a \ref{fig:ui_monitoreo_seguimientos} & RF-43 a RF-47 \\ \hline
        Perfil y ayuda & \ref{fig:ui_perfil} y \ref{fig:ui_perfil_cuenta} & RF-09, RF-12 a RF-14, RF-42, RF-48 \\ \hline
    \end{tabular}
\end{table}

\begin{figure}[htbp]
    \centering
    \pantallaUI{modo_claro/sesion_inicio.jpg}{Inicio de sesión}\hspace{0.04\linewidth}%
    \pantallaUI{modo_claro/sesion_registro.jpg}{Registro}\hspace{0.04\linewidth}%
    \pantallaUI{modo_claro/sesion_terminos.jpg}{Términos y aviso de privacidad}

    \bigskip
    \pantallaUI{modo_claro/sesion_terminos_aceptados.jpg}{Términos aceptados}\hspace{0.04\linewidth}%
    \pantallaUI{modo_claro/sesion_registro_aceptado.jpg}{Registro con los términos aceptados}
    \caption{Inicio de sesión y registro}
    \label{fig:ui_sesion_registro}
\end{figure}

\begin{figure}[htbp]
    \centering
    \pantallaUI{modo_claro/sesion_verificar_correo.jpg}{Verificación pendiente}\hspace{0.04\linewidth}%
    \pantallaUI{modo_claro/sesion_acceso_restringido.jpg}{Acceso restringido sin verificación}\hspace{0.04\linewidth}%
    \pantallaUI{modo_claro/sesion_verificacion_expirada.jpg}{Enlace de verificación no válido}
    \caption{Verificación del correo electrónico}
    \label{fig:ui_sesion_verificacion}
\end{figure}

\begin{figure}[htbp]
    \centering
    \pantallaUI{modo_claro/sesion_recuperar.jpg}{Recuperar contraseña}\hspace{0.04\linewidth}%
    \pantallaUI{modo_claro/sesion_recuperar_enviado.jpg}{Enlace de recuperación enviado}\hspace{0.04\linewidth}%
    \pantallaUI{modo_claro/sesion_nueva_contrasena.jpg}{Nueva contraseña}

    \bigskip
    \pantallaUI{modo_claro/sesion_contrasena_actualizada.jpg}{Contraseña actualizada}\hspace{0.04\linewidth}%
    \pantallaUI{modo_claro/sesion_recuperacion_expirada.jpg}{Enlace de recuperación no válido}
    \caption{Recuperación de contraseña}
    \label{fig:ui_sesion_recuperacion}
\end{figure}

\begin{figure}[htbp]
    \centering
    \pantallaUI{modo_claro/mapas_permiso_ubicacion.jpg}{Permiso de ubicación}\hspace{0.04\linewidth}%
    \pantallaUI{modo_claro/mapas_principal.jpg}{Pantalla principal}\hspace{0.04\linewidth}%
    \pantallaUI{modo_claro/mapas_busqueda_recientes.jpg}{Búsquedas recientes}

    \bigskip
    \pantallaUI{modo_claro/mapas_busqueda_sin_recientes.jpg}{Sin búsquedas recientes}\hspace{0.04\linewidth}%
    \pantallaUI{modo_claro/mapas_busqueda_resultados.jpg}{Resultados de búsqueda}\hspace{0.04\linewidth}%
    \pantallaUI{modo_claro/mapas_busqueda_sin_resultados.jpg}{Búsqueda sin coincidencias}
    \caption{Pantalla principal y búsqueda de lugares}
    \label{fig:ui_mapa_busqueda}
\end{figure}

\begin{figure}[htbp]
    \centering
    \pantallaUI{modo_claro/mapas_lugar_buscado.jpg}{Lugar buscado}\hspace{0.04\linewidth}%
    \pantallaUI{modo_claro/mapas_seleccion_lugar.jpg}{Selección de lugar}\hspace{0.04\linewidth}%
    \pantallaUI{modo_claro/mapas_seleccion_favorito.jpg}{Lugar guardado en favoritos}
    \caption{Detalle y selección de un lugar}
    \label{fig:ui_mapa_lugar}
\end{figure}

\begin{figure}[htbp]
    \centering
    \pantallaUI{modo_claro/rutas_optima.jpg}{Ruta óptima}\hspace{0.04\linewidth}%
    \pantallaUI{modo_claro/rutas_segura.jpg}{Ruta segura}\hspace{0.04\linewidth}%
    \pantallaUI{modo_claro/rutas_rapida.jpg}{Ruta rápida}

    \bigskip
    \pantallaUI{modo_claro/rutas_coincidentes.jpg}{Criterios con la misma ruta}\hspace{0.04\linewidth}%
    \pantallaUI{modo_claro/rutas_detalle.jpg}{Detalle de la ruta}\hspace{0.04\linewidth}%
    \pantallaUI{modo_claro/rutas_detalle_expandido.jpg}{Resumen de la ruta expandido}
    \caption{Selección de ruta}
    \label{fig:ui_rutas}
\end{figure}

\begin{figure}[htbp]
    \centering
    \pantallaUI{modo_claro/rutas_franja_ahora.jpg}{Franja de la hora actual}\hspace{0.04\linewidth}%
    \pantallaUI{modo_claro/rutas_franja_diurna.jpg}{Franja diurna}\hspace{0.04\linewidth}%
    \pantallaUI{modo_claro/rutas_franja_nocturna.jpg}{Franja nocturna}

    \bigskip
    \pantallaUI{modo_claro/rutas_inexistente.jpg}{Ruta peatonal no encontrada}
    \caption{Franja horaria y ruta inexistente}
    \label{fig:ui_rutas_franja}
\end{figure}

\begin{figure}[htbp]
    \centering
    \pantallaUI{modo_claro/navegacion.jpg}{Navegación}\hspace{0.04\linewidth}%
    \pantallaUI{modo_claro/navegacion_compartir.jpg}{Compartir la ruta}\hspace{0.04\linewidth}%
    \pantallaUI{modo_claro/navegacion_desvio.jpg}{Desvío detectado}

    \bigskip
    \pantallaUI{modo_claro/navegacion_recalculo.jpg}{Nueva ruta trazada}\hspace{0.04\linewidth}%
    \pantallaUI{modo_claro/navegacion_recalculo_fallido.jpg}{Recálculo fallido}\hspace{0.04\linewidth}%
    \pantallaUI{modo_claro/navegacion_sin_internet.jpg}{Sin conexión a internet}
    \caption{Navegación y recálculo de la ruta}
    \label{fig:ui_navegacion}
\end{figure}

\begin{figure}[htbp]
    \centering
    \pantallaUI{modo_claro/navegacion_emergencia.jpg}{Llamada a emergencias}\hspace{0.04\linewidth}%
    \pantallaUI{modo_claro/navegacion_monitoreo_activo.jpg}{Aviso de monitoreo activo}\hspace{0.04\linewidth}%
    \pantallaUI{modo_claro/navegacion_calificar.jpg}{Calificación del trayecto}
    \caption{Emergencias y fin del trayecto}
    \label{fig:ui_navegacion_fin}
\end{figure}

\begin{figure}[htbp]
    \centering
    \pantallaUI{modo_claro/reportes_nuevo.jpg}{Reporte desde el mapa}\hspace{0.04\linewidth}%
    \pantallaUI{modo_claro/reportes_nuevo_navegacion.jpg}{Reporte durante la navegación}\hspace{0.04\linewidth}%
    \pantallaUI{modo_claro/reportes_enviado.jpg}{Reporte enviado}

    \bigskip
    \pantallaUI{modo_claro/reportes_creado.jpg}{Reporte creado}\hspace{0.04\linewidth}%
    \pantallaUI{modo_claro/reportes_limite.jpg}{Límite de reportes alcanzado}\hspace{0.04\linewidth}%
    \pantallaUI{modo_claro/reportes_aviso_cercano.jpg}{Aviso de reporte cercano}
    \caption{Registro de reportes}
    \label{fig:ui_reportes_registro}
\end{figure}

\begin{figure}[htbp]
    \centering
    \pantallaUI{modo_claro/reportes_cercanos.jpg}{Reportes cercanos}\hspace{0.04\linewidth}%
    \pantallaUI{modo_claro/reportes_detalle.jpg}{Detalle del reporte}\hspace{0.04\linewidth}%
    \pantallaUI{modo_claro/reportes_verificar.jpg}{Verificación de un reporte cercano}

    \bigskip
    \pantallaUI{modo_claro/reportes_verificar_navegacion.jpg}{Verificación durante la navegación}
    \caption{Validación de reportes}
    \label{fig:ui_reportes_validacion}
\end{figure}

\begin{figure}[htbp]
    \centering
    \pantallaUI{modo_claro/reportes_aprobado.jpg}{Reporte aprobado}\hspace{0.04\linewidth}%
    \pantallaUI{modo_claro/reportes_aprobado_2.jpg}{Reporte aprobado, segundo caso}\hspace{0.04\linewidth}%
    \pantallaUI{modo_claro/reportes_desmentido.jpg}{Reporte desmentido}

    \bigskip
    \pantallaUI{modo_claro/reportes_desmentido_2.jpg}{Reporte desmentido, segundo caso}
    \caption{Resultado de la validación de reportes}
    \label{fig:ui_reportes_resultado}
\end{figure}

\begin{figure}[htbp]
    \centering
    \pantallaUI{modo_claro/favoritos_vacio.jpg}{Sin lugares guardados}\hspace{0.04\linewidth}%
    \pantallaUI{modo_claro/favoritos_lista.jpg}{Lugares guardados}\hspace{0.04\linewidth}%
    \pantallaUI{modo_claro/favoritos_lista_reducida.jpg}{Lista contraída}

    \bigskip
    \pantallaUI{modo_claro/favoritos_anadir.jpg}{Añadir lugar}\hspace{0.04\linewidth}%
    \pantallaUI{modo_claro/favoritos_busqueda.jpg}{Búsquedas recientes}\hspace{0.04\linewidth}%
    \pantallaUI{modo_claro/favoritos_busqueda_resultados.jpg}{Resultados de búsqueda}
    \caption{Lugares favoritos}
    \label{fig:ui_favoritos}
\end{figure}

\begin{figure}[htbp]
    \centering
    \pantallaUI{modo_claro/favoritos_detalle.jpg}{Detalle del lugar}\hspace{0.04\linewidth}%
    \pantallaUI{modo_claro/favoritos_detalle_marcado.jpg}{Lugar marcado}\hspace{0.04\linewidth}%
    \pantallaUI{modo_claro/favoritos_nuevo.jpg}{Nuevo lugar}

    \bigskip
    \pantallaUI{modo_claro/favoritos_guardar.jpg}{Guardar lugar}\hspace{0.04\linewidth}%
    \pantallaUI{modo_claro/favoritos_editar.jpg}{Editar lugar}\hspace{0.04\linewidth}%
    \pantallaUI{modo_claro/favoritos_eliminar.jpg}{Eliminar de favoritos}
    \caption{Guardado y edición de lugares favoritos}
    \label{fig:ui_favoritos_edicion}
\end{figure}

\begin{figure}[htbp]
    \centering
    \pantallaUI{modo_claro/monitoreo_enlaces.jpg}{Mis enlaces}\hspace{0.04\linewidth}%
    \pantallaUI{modo_claro/monitoreo_enlace_agregar.jpg}{Agregar enlace}\hspace{0.04\linewidth}%
    \pantallaUI{modo_claro/monitoreo_enlace_editar.jpg}{Editar enlace}

    \bigskip
    \pantallaUI{modo_claro/monitoreo_enlace_compartir.jpg}{Compartir enlace}\hspace{0.04\linewidth}%
    \pantallaUI{modo_claro/monitoreo_enlace_copiado.jpg}{Código copiado}\hspace{0.04\linewidth}%
    \pantallaUI{modo_claro/monitoreo_enlace_oculto.jpg}{Enlace oculto}
    \caption{Enlaces y códigos de seguimiento}
    \label{fig:ui_monitoreo_enlaces}
\end{figure}

\begin{figure}[htbp]
    \centering
    \pantallaUI{modo_claro/monitoreo_seguidores.jpg}{Seguidores del enlace}\hspace{0.04\linewidth}%
    \pantallaUI{modo_claro/monitoreo_solicitud.jpg}{Solicitud pendiente}\hspace{0.04\linewidth}%
    \pantallaUI{modo_claro/monitoreo_solicitud_aceptada.jpg}{Solicitud aceptada}

    \bigskip
    \pantallaUI{modo_claro/monitoreo_seguidores_verificados.jpg}{Seguidores verificados}\hspace{0.04\linewidth}%
    \pantallaUI{modo_claro/monitoreo_seguidor_eliminado.jpg}{Seguidor eliminado}\hspace{0.04\linewidth}%
    \pantallaUI{modo_claro/monitoreo_enlace_eliminar.jpg}{Eliminar enlace}
    \caption{Seguidores de un enlace}
    \label{fig:ui_monitoreo_seguidores}
\end{figure}

\begin{figure}[htbp]
    \centering
    \pantallaUI{modo_claro/monitoreo_seguimientos.jpg}{Mis seguimientos}\hspace{0.04\linewidth}%
    \pantallaUI{modo_claro/monitoreo_seguimiento_agregar.jpg}{Agregar seguimiento}\hspace{0.04\linewidth}%
    \pantallaUI{modo_claro/monitoreo_seguimiento_eliminar.jpg}{Eliminar seguimiento}

    \bigskip
    \pantallaUI{modo_claro/monitoreo_ubicacion.jpg}{Ubicación de la persona seguida}\hspace{0.04\linewidth}%
    \pantallaUI{modo_claro/monitoreo_navegacion.jpg}{Navegación de la persona seguida}\hspace{0.04\linewidth}%
    \pantallaUI{modo_claro/monitoreo_emergencia.jpg}{Llamada a emergencias}
    \caption{Seguimientos y monitoreo de otra persona}
    \label{fig:ui_monitoreo_seguimientos}
\end{figure}

\begin{figure}[htbp]
    \centering
    \pantallaUI{modo_claro/perfil.jpg}{Perfil}\hspace{0.04\linewidth}%
    \pantallaUI{modo_claro/perfil_notificaciones.jpg}{Notificaciones activadas}\hspace{0.04\linewidth}%
    \pantallaUI{modo_claro/perfil_permiso_notificaciones.jpg}{Permiso de notificaciones}

    \bigskip
    \pantallaUI{modo_claro/perfil_editar.jpg}{Editar perfil}\hspace{0.04\linewidth}%
    \pantallaUI{modo_claro/perfil_cambiar_contrasena.jpg}{Cambiar contraseña}\hspace{0.04\linewidth}%
    \pantallaUI{modo_claro/perfil_contrasena_actualizada.jpg}{Contraseña actualizada}
    \caption{Perfil del usuario}
    \label{fig:ui_perfil}
\end{figure}

\begin{figure}[htbp]
    \centering
    \pantallaUI{modo_claro/perfil_eliminar_cuenta.jpg}{Eliminar cuenta}\hspace{0.04\linewidth}%
    \pantallaUI{modo_claro/perfil_aviso_privacidad.jpg}{Aviso de privacidad}\hspace{0.04\linewidth}%
    \pantallaUI{modo_claro/perfil_manuales.jpg}{Manual de usuario}
    \caption{Cuenta y ayuda}
    \label{fig:ui_perfil_cuenta}
\end{figure}

\subsection{Modo oscuro}
El modo oscuro, que el usuario elige en la sección Apariencia de su perfil (figura \ref{fig:ui_perfil}), conserva la estructura y los flujos del modo claro y solo cambia los valores de la paleta (tabla \ref{tab:paleta_colores}): fondos gris oscuro y encabezados en café, con los botones de acción en naranja terracota, mientras que el mapa conserva sus colores. Las figuras \ref{fig:ui_oscuro_1} y \ref{fig:ui_oscuro_2} muestran una selección de pantallas de cada flujo.

\begin{figure}[htbp]
    \centering
    \pantallaUI{modo_oscuro/sesion_inicio.jpg}{Inicio de sesión}\hspace{0.04\linewidth}%
    \pantallaUI{modo_oscuro/sesion_registro.jpg}{Registro}\hspace{0.04\linewidth}%
    \pantallaUI{modo_oscuro/mapas_principal.jpg}{Pantalla principal}

    \bigskip
    \pantallaUI{modo_oscuro/rutas_optima.jpg}{Ruta óptima}\hspace{0.04\linewidth}%
    \pantallaUI{modo_oscuro/navegacion.jpg}{Navegación}\hspace{0.04\linewidth}%
    \pantallaUI{modo_oscuro/reportes_nuevo.jpg}{Reporte desde el mapa}
    \caption{Selección de pantallas en modo oscuro: sesión, mapa, rutas y reportes}
    \label{fig:ui_oscuro_1}
\end{figure}

\begin{figure}[htbp]
    \centering
    \pantallaUI{modo_oscuro/favoritos_lista.jpg}{Lugares guardados}\hspace{0.04\linewidth}%
    \pantallaUI{modo_oscuro/monitoreo_enlaces.jpg}{Mis enlaces}\hspace{0.04\linewidth}%
    \pantallaUI{modo_oscuro/monitoreo_ubicacion.jpg}{Ubicación de la persona seguida}

    \bigskip
    \pantallaUI{modo_oscuro/perfil.jpg}{Perfil}\hspace{0.04\linewidth}%
    \pantallaUI{modo_oscuro/rutas_franja_ahora.jpg}{Franja de la hora actual}\hspace{0.04\linewidth}%
    \pantallaUI{modo_oscuro/navegacion_emergencia.jpg}{Llamada a emergencias}
    \caption{Selección de pantallas en modo oscuro: favoritos, monitoreo, perfil y alertas}
    \label{fig:ui_oscuro_2}
\end{figure}
